\chapter{选学专题：高斯过程与预测不确定性}
高斯过程回归（Gaussian process regression, GPR）是一种对函数进行概率建模的方法。在给定核函数及其参数和噪声模型后，可以对输入位置形成预测均值与模型假设条件下的不确定性估计。它适合某些中小样本的平滑函数建模，能够连接物理学生熟悉的插值与统计建模。
\section{从相似输入到函数相关性}
这里的\textbf{先验}描述看当前训练数据之前对函数形状的假设，\textbf{后验}则是在读入数据后对这些可能函数的更新。高斯过程规定：在任意有限组输入位置取出的函数值，服从一个多元高斯分布；均值函数描述平均水平，协方差函数描述不同位置的函数值怎样共同变化。\textbf{协方差}可视为对两量共同偏离各自均值的平均程度的描述，与第13章输入协方差的思想相通，但此处对象是函数值。

预测均值是模型给出的中心估计，预测标准差描述给定模型假设下的分散程度。比如均值2、标准差0.3，不等于“任何真实测量都必然落在1.7到2.3内”；要把它转成区间还需要相应分布假设，并检查实际覆盖情况。

核函数 $k(\mathbf x,\mathbf z)$ 描述不同输入位置对应函数值之间的先验相关性。常见 RBF 核随输入距离增加而衰减。选择核函数相当于引入对函数平滑程度、尺度和变化方式的假设。与第11章的核 SVM 不同，GPR 的输出包含与概率模型有关的预测方差，但这种方差必须连同模型假设解释。
\section{最小 sklearn 实验}
\begin{lstlisting}[caption={人工生成平滑函数的高斯过程回归示例}]
import numpy as np
from sklearn.gaussian_process import GaussianProcessRegressor
from sklearn.gaussian_process.kernels import ConstantKernel, RBF, WhiteKernel
rng = np.random.default_rng(42)
X = np.linspace(0, 6, 30)[:, None]
y = np.sin(X[:, 0]) + rng.normal(0, 0.1, len(X))
kernel = ConstantKernel(1.0) * RBF(1.0) + WhiteKernel(0.02)
reg = GaussianProcessRegressor(kernel=kernel, normalize_y=True,
                               random_state=42).fit(X, y)
xx = np.linspace(0, 6, 200)[:, None]
mean, std = reg.predict(xx, return_std=True)
\end{lstlisting}
这里的样本由程序生成，不是天文或物理观测。实际项目可尝试小规模光变曲线或实验响应曲线，但应检查测量噪声、时间间隔、核假设及计算成本。核矩阵方法在样本数量较大时可能出现显著的存储与计算负担。
\section{不同不确定性的区别}
GPR 的预测标准差不是自动覆盖所有不确定性的万能数值。测量误差、核函数假设偏差、训练样本稀缺、分布变化和目标本身的随机性都应分别分析。学生应在图中标明区间是依据什么模型构造的，以及它是否得到独立数据的覆盖率验证。


\section{从函数集合的概率描述开始}
高斯过程可以理解为对未知函数的一组联合概率假设：在任意有限输入位置上，函数值向量服从多元高斯分布。均值函数说明先验中心，核函数说明两个输入处函数值的相关结构。这里的“先验”是模型假设，不等于已经知道真实函数；数据会通过条件分布改变预测。

设训练观测满足 $\mathbf y=\mathbf f+\boldsymbol\epsilon$，独立高斯噪声方差为 $\sigma_n^2$，函数先验均值为0、协方差矩阵为 $K$。对新输入 $x_*$，与训练位置的协方差向量为 $\mathbf k_*$。固定核与噪声参数后，条件均值和潜在函数方差为
\[
\mu_* = \mathbf k_*^\mathsf T(K+\sigma_n^2I)^{-1}\mathbf y,
\qquad
v_* = k(x_*,x_*)-\mathbf k_*^\mathsf T(K+\sigma_n^2I)^{-1}\mathbf k_*.
\]
实际计算应通过线性方程求解或Cholesky分解，不显式求逆。若预测下一次带噪观测，还需在适当独立假设下加上观测噪声方差。核参数若由数据估计，简单条件方差通常没有完整积分其估计不确定性。

\section{两个观测点的一次完整计算}
取输入0、1，观测0、1，RBF核 $k(x,z)=\exp[-(x-z)^2/2]$，噪声标准差0.1。则
\[
K+\sigma_n^2I=\begin{pmatrix}1.01&e^{-1/2}\\e^{-1/2}&1.01\end{pmatrix},
\quad\mathbf k_*=(e^{-1/8},e^{-1/8})^\mathsf T
\]
对应新输入0.5。数值求解得到预测均值约0.5459，潜在函数方差约0.03645；预测一个带同等独立噪声的新观测时，方差还要增加0.01。均值没有恰好等于两点算术平均，是因为零均值先验、相关结构与噪声共同决定了条件估计。

这个例子不是要学生相信所有函数都应服从该核。若真实响应有突变，过于平滑的核可能产生过度平滑甚至过度自信的结果。应改变核或噪声假设检查敏感性，并用独立数据检验预测覆盖，而不是仅因为程序输出了标准差就宣称不确定性完整。

\section{与本书其他方法的关系}
核SVM利用核定义相似性以形成判别函数，高斯过程则通过核定义函数值的协方差并形成条件分布。两者共享某些数学工具，目标与输出解释不同。随机森林树间离散、Bootstrap预测变化、GP条件方差也不能直接当成同一种不确定性。

高斯过程适合小到中等规模函数建模的教学，但标准稠密矩阵方法成本随样本增长较快。应用于光变曲线时，还要处理观测时刻、噪声与周期或长期趋势核；应用于物理响应时，要检查单位和领域约束。它是本书向概率建模延伸的入口，不是绕过数据与验证工作的捷径。

